\documentclass[11pt,aps,preprint,nofootinbib,superscriptaddress,tightenlines]{revtex4-2}
\usepackage[T1]{fontenc}
\usepackage{amsmath,amssymb,mathtools,bm}
\usepackage{xcolor}
\usepackage{booktabs}
\usepackage{graphicx}
\usepackage{dsfont}
\newcommand{\one}{\mathds{1}}
\newcommand{\avg}[1]{\langle #1\rangle}
\definecolor{darkblue}{RGB}{20,60,140}
\usepackage[colorlinks=true,linkcolor=darkblue,citecolor=darkblue,urlcolor=darkblue]{hyperref}

\def\st{\begin{equation}}
\def\stp{\end{equation}}
\def\Eq#1{eq.~(\ref{#1})}
\def\Sect#1{Sect.~\ref{#1}}
\def\Fig#1{Fig.~\ref{#1}}
\let\oldparagraph\paragraph
\renewcommand{\paragraph}{\par\bigskip\oldparagraph}

\begin{document}
\sloppy

\title{The dynamics with the split Hamiltonian: campaign \texttt{hastings\_split}}
\author{D.~Teaney}
\affiliation{Department of Physics and Astronomy, Stony Brook University,
Stony Brook, New York 11794, USA}
\date{\today}

\begin{abstract}
The split renormalization (Step 5, \texttt{docs/step5\_plan.tex}) was tested so far in equilibrium only, by hybrid Monte Carlo
(campaign \texttt{eos\_split}).  Here the full dynamics --- the Hamiltonian step and the stochastic viscous step \texttt{moment} with the
Hastings correction, the current defaults --- runs with the split Hamiltonian (\texttt{SPH\_RENORM=split}) and, with the same seeds and
settings, at the tree level (no renormalization).  The questions: (i) does the dynamics with the split Hamiltonian sample its own
equilibrium (dynamics against HMC of the same Hamiltonian)? (ii) how much do the equal-time spectra and the standard plots of
\texttt{docs/inverse\_moment.tex} (shear damping, sound, the correlators along and across the flow at $v_0=0.8$) change when the
equation of state is renormalized?  The intention is to make \texttt{SPH\_RENORM=split} the default once these answers are in.
\end{abstract}

\maketitle

%=============================================================================
\section{What was run}
\label{sec:runs}

\paragraph{Provenance rule.}  Every figure caption in this note begins with the campaign, the run sets it reads (the directory names
under \texttt{\$SCRATCH/sph/<campaign>/}) and the commit; the table below gives each campaign's settings.  Every run also records its
settings (\texttt{cs2\_val}, \texttt{SPH\_RENORM}, the commit) in its \texttt{.toml}.

\paragraph{Campaign \texttt{hastings\_split}} (\texttt{nersc/campaign\_hastings\_split.txt}).  A uniform two-dimensional conformal fluid
in $d=2$, $c_s^2=1/2$ (the default equation of state of \texttt{examples/equilibrium.jl} and \texttt{examples/hmc\_positions.jl}),
$\eta/s=0.2$, periodic box $L=0.6$ ($N=672$), $h=0.05$, $h/a=2$, energy Hamiltonian with the number volume, generalized leapfrog,
Liouville measure, HMC start, viscous step \texttt{moment} with the Hastings correction, $\gamma_{\max}=100$ (all current defaults).
Entropy per particle $s_a=16$ ($S_0=2.956\times10^4$) and $54$ ($9.977\times10^4$); fluid at rest and moving with $v_0=0.8$ along $x$.
\begin{center}
\begin{tabular}{llll}
\toprule
run sets & Hamiltonian & what & per set \\ \midrule
\texttt{hs<sa>\_tr\_v<v>} & tree level (\texttt{SPH\_RENORM=0}) & full dynamics, $t=100$ & 16 seeds \\
\texttt{hs<sa>\_sp\_v<v>} & split (\texttt{SPH\_RENORM=split}) & full dynamics, $t=100$ & 16 seeds (same as tr) \\
\texttt{hq<sa>\_tr\_v<v>} & tree level & HMC, 300 updates & 16 chains \\
\texttt{hq<sa>\_sp\_v<v>} & split & HMC, 300 updates & 16 chains \\
\bottomrule
\end{tabular}
\end{center}
$\langle sa\rangle\in\{16,54\}$, $\langle v\rangle\in\{0,0.8\}$: 256 runs.  Time-step factor 0.6 at $v_0=0.8$; HMC trajectories of random
length in $[\tau/2,\tau]$, $\tau=L/c_s$, $\delta t$ factor 0.35.

\paragraph{The earlier campaign \texttt{hastings} is not a reference.}  Campaign \texttt{hastings} and its reruns
(\texttt{hastings\_s16}, \texttt{hastings\_hq16}, \texttt{hastings\_low}; \texttt{inverse\_moment.tex}, Stage 2, figures
\texttt{hastings\_nersc\_*}) ran at commits \texttt{927e035}, \texttt{e282211}, \texttt{77cf74f}, before the default was changed to
$d=d_s$ (\texttt{8d50d04}): their fluid has $c_s^2=1/3$ (physical dimension 3 in a two-dimensional simulation), recorded as
\texttt{cs2\_val = 0.333} in every run.  The present tree-level lines repeat their \texttt{h} and \texttt{hq} lines at $c_s^2=1/2$.
Nothing else differs between the tr and sp lines of \texttt{hastings\_split}: the same seeds, settings and code; only
\texttt{SPH\_RENORM}.

\paragraph{What the split Hamiltonian changes at $c_s^2=1/2$.}  For the conformal fluid in $d=2$, $\mathcal N=1-2c_s^2=0$: the counterterm
is the label part $b\ln u$ with $b=5/2$ in the measure and the motion term $E_{\rm dyn}=\tfrac12T\ln(1-c_s^2v^2)$ in $H$
(\texttt{step5\_plan.tex}).  The equilibrium effect (campaign \texttt{eos\_split}, Figs.~1--6 of \texttt{step5\_plan.tex}) is to remove the
$O(1/s)$ error of the tree level in $E$, $P$ and the susceptibilities.  The equal-time spectra at small $k$ are the susceptibilities,
so a change of order $1/s_a$ (a few percent at $s_a=16$) is expected there.

%=============================================================================
\section{Results}
\label{sec:results}

\emph{Pending: the campaign is running.  Planned figures (each with its campaign and run sets in the caption):}
\begin{enumerate}
\item Equal-time spectra $\avg{|M_{L,\bm k}|^2}$, $\avg{|M_{T,\bm k}|^2}$: tree dynamics / tree HMC, split dynamics / split HMC, split HMC /
  tree HMC (\texttt{HA\_SPLIT=1 examples/hastings\_analysis.jl}; \texttt{hastings\_split\_spectra.png}).
\item For $s_a=54$ and $16$ (\texttt{BOOST\_GROUP=SP54}, \texttt{SP16} \texttt{examples/boost\_analysis.jl}): the equal-time correlators,
  the fitted damping and frequencies against hydrodynamics, and the time correlators (shear and sound at rest; shear, sound, $EE$ along the
  flow and sound, shear across it at $v_0=0.8$), tree and split overlaid.
\end{enumerate}

\end{document}
